\section{Introduction}

Ant colony optimization (ACO)~\cite{dorigo1996ant,stutzle2000maxmin} and
Physarum polycephalum conductivity
models~\cite{nakagaki2000maze,tero2007mathematical} maintain internal state
--- a pheromone field, a field of tube conductivities --- that is claimed to
let them adapt to a changed network rather than recompute from scratch. On a
road network, where closures and congestion change edge costs continuously,
that is an attractive property.

It is also rarely measured. Evaluations of these methods overwhelmingly use
synthetic graphs that are undirected and comparatively dense, report accuracy
on an undisrupted shortest-path problem, and treat the adaptive machinery as
self-evidently useful. Three gaps follow: the test graphs lack the structure
that makes real road routing hard (one-way streets, low out-degree,
cul-de-sacs); congestion is usually absent, so the weights the adaptation
tracks never move; and recovery after disruption --- the behaviour the
retained state exists to provide --- is seldom the measured quantity.

We close those gaps for one city. We benchmark Dijkstra and
A*~\cite{dijkstra1959note,hart1968formal}, a MAX--MIN Ant System, and a
Tero-style Physarum model on the real road network of central Bengaluru
under congestion from the SUMO microsimulator~\cite{lopez2018microscopic},
measuring recovery from seven classes of closure under warm start (retained
state) and cold start (state reset). Every algorithm sees identical OD
pairs, closures and traffic, and all quality figures use a Dijkstra optimum
computed on the same weighted graph.

\textbf{Contribution 1: extraction provenance, not the model, decides what a
Physarum result means.} The conductivity model emits a field, not a route; a
separate extraction step produces the path. Instrumenting that step shows
Physarum's apparent competence is largely its fallback's. Pooled, it matches
the optimum on 90.0\% of undisrupted pairs, reproducing the literature;
restricted to its own maximum-flux walk, 23.3\%. The median stretch is
exactly 1.000 in \emph{every} provenance class, so a pooled number cannot
distinguish the model from its Dijkstra fallback. On warm recoveries the
fallback supplies 66.8\% of reported paths.

\textbf{Contribution 2: warm start buys speed and charges quality.} It cuts
recovery 2.97$\times$ in iterations and 6.87$\times$ in wall-clock, paired
($n=358$); 48.3\% of warm recoveries converge in exactly one iteration,
which cold start never does. Per-iteration cost is unchanged, so the saving
is adaptation rather than cheaper arithmetic. Against this, warm start is
significantly worse on the reported post-disruption path
($p=4.3\times10^{-7}$) and completes its own extraction far less often.

We report what the methods lose as well as what they gain: Dijkstra
recomputes an exact route in one iteration and about 31\,ms, and no method
here approaches it on wall-clock.
